\DocumentMetadata{
  pdfversion=2.0,
  pdfstandard=ua-2,
  lang=en-US,
  tagging=on,
  tagging-setup={math/setup=mathml-SE,tabsorder=structure}
}
\documentclass[11pt,letterpaper]{article}
\usepackage[margin=0.68in,includefoot]{geometry}
\usepackage{newcomputermodern}
\usepackage{amsmath,amscd,mathtools}
\providecommand{\square}{\mdlgwhtsquare}
\usepackage[hidelinks]{hyperref}
\hypersetup{
  pdftitle={Numerical Analysis PhD Exam, January 2020},
  pdfauthor={University of Florida Department of Mathematics},
  pdfsubject={Graduate mathematics examination},
  pdfdisplaydoctitle=true,
  bookmarksopen=true
}
\setcounter{secnumdepth}{0}
\setlength{\parindent}{0pt}
\setlength{\parskip}{0.28em}
\setlength{\emergencystretch}{3em}
\setlength{\leftmargini}{2.15em}
\setlength{\leftmarginii}{2.65em}
\setlength{\leftmarginiii}{3em}
\renewcommand{\labelenumi}{\textbf{\theenumi.}}
\pagestyle{empty}
\raggedbottom
\allowdisplaybreaks
\newenvironment{examproblems}[1]{%
  \begin{enumerate}%
  \setcounter{enumi}{#1}%
  \setlength{\topsep}{0.35em}%
  \setlength{\partopsep}{0pt}%
  \setlength{\itemsep}{0.48em}%
  \setlength{\parsep}{0.18em}%
}{\end{enumerate}}
\newenvironment{examsubparts}{%
  \begin{enumerate}%
  \setlength{\topsep}{0.18em}%
  \setlength{\partopsep}{0pt}%
  \setlength{\itemsep}{0.16em}%
  \setlength{\parsep}{0.1em}%
}{\end{enumerate}}
\newenvironment{examinstructions}{%
  \par\begingroup\itshape
}{%
  \par\endgroup\vspace{0.18em}
}
\makeatletter
\renewcommand\section{\@startsection{section}{1}{0pt}%
  {-1.25ex plus -.25ex minus -.1ex}%
  {0.55ex plus .1ex}%
  {\normalfont\large\bfseries}}
\renewcommand{\@maketitle}{%
  \newpage\null\vskip -1.2em%
  {\footnotesize\bfseries
    \href{https://www.ufl.edu/}{University of Florida}, \href{https://math.ufl.edu/}{Department of Mathematics}\par}%
  \vskip 0.18em%
  \tagstructbegin{tag=Title}%
    {\LARGE\bfseries\href{https://gma.math.ufl.edu/exams/numerical-analysis-phd/}{Numerical Analysis PhD Exam}, January 2020\par}%
  \tagstructend%
  \vskip 0.35em%
  \tagmcbegin{artifact}%
    \hrule height 1pt%
  \tagmcend%
  \vskip 0.55em%
}
\makeatother
\title{Numerical Analysis PhD Exam, January 2020}
\author{}
\date{}
\begin{document}
\maketitle
\thispagestyle{empty}
\begin{examinstructions}
Do 4 (four) of the first 5 (1–5) and 4 (four) of the last 5 problems (6–10).
\end{examinstructions}
\begin{examproblems}{0}
\item
Let \(A \in \mathbb{C}^{m\times n}\).
\begin{examsubparts}
\item[\textnormal{(a)}]
Determine constants~\(\alpha\) and~\(\beta\) such that the following inequality holds for the~\(p\) and~\(\infty\) norms of matrix~\(A\), for integers \(p\geq 1\). Justify your answer. 
\[
\alpha\|A\|_\infty\leq \|A\|_p\leq \beta\|A\|_\infty.
\]
\item[\textnormal{(b)}]
Prove or give a counterexample: \(\|A\|_2\leq \|A\|_F\), where \(\|A\|_F\) is the Frobenius norm of~\(A\). If you prove this, make sure to justify each nontrivial step.
\end{examsubparts}
\item
Suppose~\(A\) is Hermitian positive definite.
\begin{examsubparts}
\item[\textnormal{(a)}]
Prove that each principal submatrix of~\(A\) is Hermitian positive definite.
\item[\textnormal{(b)}]
Prove that an element of~\(A\) with largest magnitude lies on the diagonal.
\item[\textnormal{(c)}]
Prove that~\(A\) has a Cholesky decomposition.
\end{examsubparts}
\item
Define the matrices~\(A\) and~\(B\) by 
\[
A=\begin{pmatrix}
1&2&0\\
1&2&0\\
1&2&0\\
1&2&0
\end{pmatrix},
\qquad
B=\begin{pmatrix}
1&0&1\\
0&2&0\\
1&0&0\\
0&0&1
\end{pmatrix}.
\]
\begin{examsubparts}
\item[\textnormal{(a)}]
Find both full and economy singular value decompositions of~\(A\).
\item[\textnormal{(b)}]
Find both full and economy QR decompositions of~\(B\).
\end{examsubparts}
\item
Let \(\|\cdot\|\) be a subordinate (induced) matrix norm.
\begin{examsubparts}
\item[\textnormal{(a)}]
If~\(E\) is \(n\times n\) with \(\|E\|<1\), then show \(I+E\) is nonsingular and 
\[
\|(I+E)^{-1}\|\leq \frac{1}{1-\|E\|}.
\]
\item[\textnormal{(b)}]
If~\(A\) is \(n\times n\) invertible and~\(E\) is \(n\times n\) with \(\|A^{-1}\|\|E\|<1\), then show \(A+E\) is nonsingular and 
\[
\|(A+E)^{-1}\|\leq \frac{\|A^{-1}\|}{1-\|A^{-1}\|\|E\|}.
\]
\end{examsubparts}
\item
Suppose the linear equation \(Ax=b\), with 
\[
A=\begin{pmatrix}\delta&1\\1&1\end{pmatrix},\qquad |\delta|<\epsilon_m/4,
\]
 is solved on a floating-point system with machine-epsilon \(\epsilon_m\), using an LU factorization (no pivoting!) followed by forward and back substitution. (You may assume the operations \((+,-,\div,\times)\), do not incur any additional errors.)
\begin{examsubparts}
\item[\textnormal{(a)}]
If \(b=(1,0)^T\), compute the backward error, \(\|\hat b-b\|_2/\|b\|_2\), where \(\hat b\) is the data that satisfies \(A\hat x=\hat b\), and \(\hat x\) is the computed solution.
\item[\textnormal{(b)}]
Is the result of (a) sufficient to draw any conclusions about the backward-stability of the algorithm used to compute \(\hat x\)? Explain.
\end{examsubparts}
\item
This problem has two parts.
\begin{examsubparts}
\item[\textnormal{(a)}]
Starting with the first two Legendre polynomials over \([-1,1]\), given by \(p_0(x)=1\), and \(p_1(x)=x\), find the next three (monic) Legendre polynomials, \(p_2(x),p_3(x)\) and \(p_4(x)\).
\item[\textnormal{(b)}]
Find the nodes \(t_0,t_1,t_2\) and weights \(w_0,w_1,w_2\) which define the Gaussian Quadrature formula 
\[
\int_{-1}^1 f(t)\,dt\approx w_0f(t_0)+w_1f(t_1)+w_2f(t_2),
\]
 with degree of exactness \(m=5\). You should find the nodes exactly, and may leave the weights \(w_0,w_1,w_2\), in integral form.
\end{examsubparts}
\item
Consider the fixed point problem \(x=f(x)\), where \(f(x)=e^{-(3+x)}\).
\begin{examsubparts}
\item[\textnormal{(a)}]
Assuming all computations are done in exact arithmetic, find the largest open interval in~\(\mathbb{R}\) where the fixed-point iteration \(x_{k+1}=f(x_k)\) is ensured to converge.
\item[\textnormal{(b)}]
Write a Newton iteration for finding the fixed-point.
\end{examsubparts}
\item
Let \(P_1\) be the space of polynomials of degree at most one. Using the norm \(\|u\|_2=\left(\int_a^b u^2\,dx\right)^{1/2}\)
\begin{examsubparts}
\item[\textnormal{(a)}]
Find the least-squares approximation to \(f(x)=x^3\) in \(P_1\) over \([a,b]=[-1,1]\).
\item[\textnormal{(b)}]
Find the least-squares approximation to \(f(x)=x^4\) in \(P_1\) over \([a,b]=[0,1]\).
\end{examsubparts}
\item
This problem has two parts.
\begin{examsubparts}
\item[\textnormal{(a)}]
Find a natural cubic spline \(B_0\) on \(-1\leq x\leq 1\) that satisfies \(v(-1)=0\), \(v(0)=2\), and \(v(1)=0\).
\item[\textnormal{(b)}]
Find a natural cubic spline on \(-1\leq x\leq 2\) that satisfies \(v(-1)=0\), \(v(0)=2\), \(v(1)=1/2\) and \(v(2)=0\). You may write the solution as \(B=B_0+B_1\), where \(B_0\) is (or is closely related to) the solution to (a), if you like.
\end{examsubparts}
\item
Let \(f\in C^\infty(a-H,a+H)\), and let \(h<H\). Let \(x_0=a-h\), \(x_1=a\) and \(x_2=a+h\).
\begin{examsubparts}
\item[\textnormal{(a)}]
Find the finite difference approximation to \(f''(a)\) based on the quadratic interpolant \(p_2\) which satisfies \(p_2(x_0)=f(x_0)\), \(p_2(x_1)=f(x_1)\) and \(p_2(x_2)=f(x_2)\) (you should explicitly show how the difference approximation is derived from the interpolant).
\item[\textnormal{(b)}]
Let \(\psi_0(h)=\psi(h)\) be the difference approximation to \(f''(a)\) found in part (a). Assume (in exact arithmetic) \(\psi(h)\to\psi(0)=f''(a)\) as \(h\to 0\), and that \(\psi(h)\) has the asymptotic expansion 
\[
\psi(h)=\psi(0)+a_2h^2+a_4h^4+a_6h^6+\cdots.
\]
 Find the general Richardson extrapolation formula for \(\psi_k(h)\) based on \(\psi_{k-1}(h)\) and \(\psi_{k-1}(h/2)\).
\end{examsubparts}
\end{examproblems}
\end{document}
